% ===================================================================
%  TABULKA č. 5 — Dům a jeho stavba.
%  Traduction 2026 d'apres texto/io, controlee sur texto/fr, texto/ru
%  et texto/uk. Consigne : voir texto/cs/10-tabulka-01.tex.
%
%  LE TABLEAU DU BATIMENT EST UN DIALOGUE, le seul du volume, et il
%  met le tcheque devant un choix que sa voisine slave n'avait pas :
%  LE TUTOIEMENT. L'ido dit « tu » a l'enfant et « vi » aux hommes ;
%  le tcheque distingue « ty » et « vy » exactement de meme, et
%  accorde en outre le participe au genre du sujet. Une phrase comme
%  « vidíš » ne se traduit donc pas sans savoir a qui elle s'adresse
%  — et le fac-simile le dit chaque fois par le nom du parleur.
%
%  LE LEXIQUE DU CHANTIER EST PRESQUE ENTIEREMENT TCHEQUE, ce qui
%  n'etait le cas dans aucune des colonnes romanes : « zedník » le
%  magon, « tesař » le charpentier, « pokrývač » le couvreur,
%  « truhlář » le menuisier, « zámečník » le serrurier, « sklenář »
%  le vitrier, « klempíř » le zingueur, « kovář » le forgeron. Huit
%  metiers, huit mots slaves, la ou le galicien en empruntait la
%  moitie a la maconnerie latine.
%
%  ET C'EST LA MEME REGLE, PRISE PAR L'AUTRE BOUT. Le chantier
%  s'apprend sur le chantier ; la ou le chantier n'a pas ete importe,
%  ses mots ne l'ont pas ete non plus.
% ===================================================================

%%K t05-apar-1 apar
\VUsaut{6.0mm}
\VUtitre{15.0pt}{60}{TABULKA č. 5}
\VUsaut{1.4mm}
\VUtitre{11.4pt}{0}{Dům a jeho stavba.}
\VUsaut{2.2mm}

%%K t05-tit-1 sub
\VUpk{10.2pt}{40}{Šťastné setkání.}
\VUsaut{3.0mm}

%%K t05-01-1 p
\textsc{Jan}. --- Strýčku, velmi bych chtěl vědět, jak se staví \VUgras{dům}.

%%K t05-01-2 p
\textsc{Strýc} \textsuperscript{(80)}. --- To je velmi snadné, příteli, pojď se mnou,
ukážu ti ten, který dává stavět pan Bernard \textsuperscript{(79)} a v němž budu
bydlet.

%%K t05-01-3 p
\textsc{Jan}. --- A kdo tedy nakreslil \VUgras{plán} \textsuperscript{(76)} tohoto
domu?

%%K t05-01-4 p
\textsc{Strýc}. --- \VUgras{Architekt} \textsuperscript{(75)}; předložil velmi jasnou
kresbu, která byla schválena, a \VUgras{stavitel} \textsuperscript{(77)} začal práce.

%%K t05-01-5 p
\textsc{Jan}. --- A kdo je ten stavitel?

%%K t05-01-6 p
\textsc{Strýc}. --- Je to pan Blanc, otec tvého přítele Jakuba. Řídí dělníky spolu s
panem Rouxem, architektem.

%%K t05-01-7 p
Nu! Zde právě včas přichází pan Blanc se svým \VUgras{pravítkem}
\textsuperscript{(78)} a pan Roux se svým plánem.

%%K t05-01-8 p
Dobrý den, pánové. Jak se máte?

%%K t05-01-9 p
\textsc{Pan Roux}. --- Velmi dobře, děkuji. Dobrý den, Jene, jak to dítě vyrostlo!

%%K t05-01-10 p
\textsc{Strýc}. --- Ano, opravdu, je to malý muž. Je velmi vážný; je třeba mu vysvětlit
všechno, co vidí. Dnes by chtěl vědět, jak se staví dům.

%%K t05-tit-2 sub
\VUpk{10.2pt}{40}{Dělníci.}
\VUsaut{3.0mm}

%%K t05-01-11 p
\textsc{Pan Blanc}. --- Nuže, pojď se mnou, kamaráde, hned ti to vysvětlím, neboť mám
velmi rád děti, které se chtějí poučit.

%%K t05-01-12 p
Vidíš toho dělníka, který drží v ruce \VUgras{lopatu} \textsuperscript{(82)}: je to
\VUgras{zemní dělník} \textsuperscript{(81)}. Kope \VUgras{rýhu}
\textsuperscript{(46)}, aby položil vodovodní potrubí. Vykopal také půdu na místě, kde
je dům, aby zedník postavil čtyři \VUgras{zdi}. Části těch zdí, které jsou v zemi, se
nazývají \VUgras{základy} \textsuperscript{(2)}. Prostor uzavřený mezi základy tvoří
\VUgras{sklep} \textsuperscript{(3)}. Ten sklep má malé okno zvané \VUgras{sklepní
okénko} \textsuperscript{(5)}, \VUgras{dveře} \textsuperscript{(4)} a schody.

%%K t05-01-13 p
\VUgras{Zedník} \textsuperscript{(108)} pokračoval ve stavbě zdí a nechával otvory pro
\VUgras{dveře} \textsuperscript{(8)} a \VUgras{okna} \textsuperscript{(17)}.

%%K t05-01-14 p
\textsc{Jan}. --- Ale toho zedníka nevidím!

%%K t05-01-15 p
\textsc{Pan Blanc}. --- Teď skončil práci v domě. Je blízko \VUgras{stáje}
\textsuperscript{(54)}, na svém \VUgras{lešení} \textsuperscript{(110)}, staví zeď
\VUgras{prádelny} \textsuperscript{(56)}.

%%K t05-01-16 p
Má \VUgras{zednickou lžíci}, \VUgras{necky} a \VUgras{olovnici}
\textsuperscript{(109)}. Ale ta jména si nezapamatuješ.

%%K t05-01-17 p
\textsc{Jan}. --- Ale ano, hned si vyžádám od strýce malou zednickou lžíci a necky, a
sám si vyrobím olovnici z provázku.

%%K t05-tit-3 sub
\VUsaut{3.0mm}
\VUpk{10.2pt}{40}{Materiál.}
\VUsaut{3.0mm}

%%K t05-01-18 p
\textsc{Strýc}. --- Ha! Velmi dobře. Ale řekni mi, Jene, čeho je třeba ke stavbě domu?

%%K t05-01-19 p
\textsc{Jan}. --- Je třeba \VUgras{kvádrů} \textsuperscript{(59)} a \VUgras{malty}
\textsuperscript{(65)}.

%%K t05-01-20 p
\textsc{Strýc}. --- Správně, podívej se na toho muže s velkým kloboukem na svém
\VUgras{voze} \textsuperscript{(136)}. Je to \VUgras{vozka} \textsuperscript{(135)}.
Každý den sem přiváží \VUgras{kamenné bloky} \textsuperscript{(60)}. Skládá svůj vůz na
dvoře a \VUgras{kameníci} \textsuperscript{(83)} tešou bloky a řežou je svou
\VUgras{kamenickou pilou} \textsuperscript{(85)}. \VUgras{Pomocný dělník}
\textsuperscript{(146)} je nosí zedníkovi, který je klade.

%%K t05-01-21 p
Zatím \VUgras{míchač} \textsuperscript{(87)} připravuje maltu. Potřebuje \VUgras{písek}
\textsuperscript{(62)}, \VUgras{vápno} \textsuperscript{(63)} a \VUgras{vodu}
\textsuperscript{(64)}. Vápno je blízko něho v \VUgras{pytli} \textsuperscript{(71)},
\VUgras{zednický přisluhovač} \textsuperscript{(111)} mu přiváží písek na
\VUgras{kolečku} \textsuperscript{(112)} a \VUgras{učedník} \textsuperscript{(113)} je
u pumpy (58). Plní \VUgras{kbelík} \textsuperscript{(114)}. Malta je brzy hotova.
\VUgras{Nosič malty} \textsuperscript{(107)} ji bere a nese nahoru po žebříku na
lešení, kde pracuje zedník, který dokončuje tu stranu domu.

%%K t05-01-22 p
A nyní, jak se nazývá dělník, který pracuje na \VUgras{střeše}
\textsuperscript{(31)}?

%%K t05-01-23 p
\textsc{Jan}. --- Je to \VUgras{pokrývač} \textsuperscript{(118)}.

%%K t05-tit-4 sub
\VUsaut{3.0mm}
\VUpk{10.2pt}{40}{Střecha domu.}
\VUsaut{3.0mm}

%%K t05-01-24 p
\textsc{Pan Blanc}. --- Ano, ale nejprve přichází \VUgras{tesař}
\textsuperscript{(115)}.

%%K t05-01-25 p
Tesař zpracovává \VUgras{krov} \textsuperscript{(35)}. Spojuje \VUgras{trámy}
\textsuperscript{(37)} \VUgras{kolíky} \textsuperscript{(117)}. Ti dělníci tam nahoře,
se svými širokými sametovými rukávy, jsou tesaři. Jeden drží trámek a druhý řeže kolík
svou \VUgras{sekerou} \textsuperscript{(116)}. Ten, který je blízko \VUgras{komína}
\textsuperscript{(41)}, je pokrývač. Přibíjí latě na (*) \VUgras{stropnice} a
\VUgras{břidlice} \textsuperscript{(66)} na latě. Někdy klade tašky.

%%K t05-noto-1 noto
\VUnotes{143.2mm}{%
(*) \VUgras{Balk-o} = něm. \textit{Balken}, angl. \textit{joist}, fr.
\textit{solive}, it. \textit{travicello}, rus. \textit{balka}, čes.
\textit{stropnice}, šp. a port. \textit{viga}. Odborný kořen dosud nepřijatý, ale
hodný návrhu k vyjádření smyslu francouzského \textit{solive}, jež není ani
\textit{trabo} (fr. \textit{poutre}), ani trámek. Je to otesaná dřevěná tyč, která
nese podlahu a strop.%
}

%%K t05-01-26 p
Střecha stáje je \VUgras{kryta taškami} \textsuperscript{(55)}, střecha domu je
\VUgras{kryta břidlicí} \textsuperscript{(32)} a střecha verandy je \VUgras{zinková}
\textsuperscript{(24)}.

%%K t05-01-27 p
\textsc{Jan}. --- A pracuje pokrývač také na zinkových střechách?

%%K t05-01-28 p
\textsc{Pan Blanc}. --- Ne, to dělá \VUgras{klempíř} \textsuperscript{(121)} nebo
\VUgras{olovník}, ten, který klade \VUgras{vikýř} \textsuperscript{(38)} na střechu
domu. Klade také \VUgras{okapové roury} \textsuperscript{(43)} a \VUgras{střešní
žlaby} \textsuperscript{(42)}. Spájí zinek a plech. Připevnil \VUgras{hromosvod}
\textsuperscript{(40)} a \VUgras{větrnou korouhev} \textsuperscript{(39)} na
\VUgras{štít} \textsuperscript{(36)}.

%%K t05-01-29 p
\textsc{Jan}. --- A tak je dům dokončen?

%%K t05-tit-5 sub
\VUsaut{3.0mm}
\VUpk{10.2pt}{40}{Nástroje.}
\VUsaut{3.0mm}

%%K t05-01-30 p
\textsc{Pan Blanc}. --- Ne zcela. \VUgras{Štukatér} \textsuperscript{(99)} staví z
\VUgras{cihel} \textsuperscript{(61)} příčky, které jej dělí na různé pokoje. Potahuje
\VUgras{sádrou} \textsuperscript{(70)} \VUgras{stropy} \textsuperscript{(26)} a zdi.
Teď pracuje na stropě \VUgras{verandy} \textsuperscript{(33)}. Má, stejně jako zedník,
\VUgras{zednickou lžíci} \textsuperscript{(100)} a \VUgras{necky}
\textsuperscript{(101)}.

%%K t05-01-31 p
Potom truhlář zpracovává dveře, okna, \VUgras{parkety} \textsuperscript{(16)} a
\VUgras{okenice} \textsuperscript{(19)}.

%%K t05-01-32 p
Teď pracuje na dvoře na svém \VUgras{pracovním stole} \textsuperscript{(90)}. Má
\VUgras{pilu} \textsuperscript{(94)}, aby řezal dřevo (69), \VUgras{hoblík}
\textsuperscript{(91)}, \VUgras{svěrku} \textsuperscript{(92)}, \VUgras{dláto}
\textsuperscript{(94 bis)}, \VUgras{kružítko} \textsuperscript{(95)} a dřevěnou
\VUgras{palici} \textsuperscript{(95 bis)}.

%%K t05-01-33 p
\textsc{Jan}. --- Ha! Ale je zábavnější dělat truhláře než zedníka. Strýčku, koupíš mi
pracovní stůl, pilu a hoblík, neboť si brzy udělám boudu pro Medora.

%%K t05-01-34 p
\textsc{Strýc}. --- Ano, chlapečku.

%%K t05-01-35 p
\textsc{Jan}. --- Jak jsem spokojen! A kdo je ten, který stojí blízko vchodových dveří?

%%K t05-tit-6 sub
\VUsaut{3.0mm}
\VUpk{10.2pt}{40}{Natírání.}
\VUsaut{3.0mm}

%%K t05-01-36 p
\textsc{Pan Blanc}. --- Je to \VUgras{natěrač} \textsuperscript{(102)}. Stojí na svém
žebříku s hrncem \VUgras{barvy} \textsuperscript{(103)} a se svým \VUgras{štětcem}
\textsuperscript{(104)}. Natírá dřevo vchodových dveří, aby se déle uchovalo.

%%K t05-01-37 p
On také lepí tapety v bytech.

%%K t05-01-38 p
\VUgras{Čalouník} \textsuperscript{(107)}, který je na \VUgras{štaflích}
\textsuperscript{(106)}, na \VUgras{balkoně} \textsuperscript{(28)}, klade
\VUgras{roletu} \textsuperscript{(29)}.

%%K t05-01-39 p
Dělník, který stojí blízko velkého okna, je \VUgras{sklenář} \textsuperscript{(96)}.
Připevňuje \VUgras{okenní tabulky} \textsuperscript{(18)} \VUgras{tmelem}
\textsuperscript{(73)}. Ten druhý dělník, klečící blízko sklepních dveří, je
\VUgras{zámečník} \textsuperscript{(97)}. Klade \VUgras{zámek} \textsuperscript{(6)}.
Jeho \VUgras{pilník} \textsuperscript{(98)} leží blízko něho.

%%K t05-01-40 p
\textsc{Jan}. --- A je také zámečníkem ten dělník, který tluče svým \VUgras{kladivem}
\textsuperscript{(84)} do železa?

%%K t05-01-41 p
\textsc{Pan Blanc}. --- Přesněji řečeno, je to kovář (125). Kove \VUgras{železné
součásti} \textsuperscript{(67)} pro \VUgras{elektrikáře} \textsuperscript{(124)},
který je na střeše \VUgras{kolny} \textsuperscript{(51)}. Má \VUgras{výheň}
\textsuperscript{(126)}, \VUgras{kovadlinu} \textsuperscript{(127)} a
\VUgras{kleště} \textsuperscript{(128)}.

%%K t05-01-42 p
Elektrikář připevňuje \VUgras{měděný} \textsuperscript{(44)} \VUgras{drát} telefonu.

%%K t05-tit-7 sub
\VUpk{10.2pt}{40}{Zahradničení.}
\VUsaut{3.0mm}

%%K t05-01-43 p
\textsc{Strýc}. --- Vzadu na \VUgras{dvoře} \textsuperscript{(45)}, blízko
\VUgras{mříže} \textsuperscript{(47)} a \VUgras{vrat} \textsuperscript{(48)}, vidíš
\VUgras{zahradníka} \textsuperscript{(129)}. Připravuje \VUgras{zahrádku}
\textsuperscript{(49)}. Zryl zemi svým \VUgras{rýčem} \textsuperscript{(131)}, seje
semena a hrabe \VUgras{hráběmi} \textsuperscript{(130)}. Potom svou \VUgras{konví}
\textsuperscript{(132)} zalije \VUgras{záhony} \textsuperscript{(150)} a rostliny
porostou.

%%K t05-01-44 p
\VUgras{Podkoní} \textsuperscript{(133)}, který právě ošetřil koně a dal mu krmivo,
čistí \VUgras{kočár} \textsuperscript{(52)} houbou, \VUgras{ruční pumpou}
\textsuperscript{(134)} a kbelíkem vody. Potom jej znovu zaveze do kolny a zavře dveře
tlustou \VUgras{závorou} \textsuperscript{(53)}.

%%K t05-01-45 p
Teď jsi spokojen, myslím, dostalo se ti dostatečných vysvětlení.

%%K t05-01-46 p
\textsc{Jan}. --- Ano, strýčku, ale rád bych viděl vnitřek domu.

%%K t05-01-47 p
\textsc{Strýc}. --- To bude zítra. Pan Roux ti vysvětlí plán a označí ti jméno každého
pokoje.

%%K t05-tit-8 sub
\VUsaut{3.0mm}
\VUpk{10.2pt}{40}{Vnitřní uspořádání domu.}
\VUsaut{2.4mm}

%%K t05-apar-2 apar
{\centering\textit{(Viz plán.)}\par}
\VUsaut{2.4mm}

%%K t05-01-48 p
\textsc{Architekt}. --- Pojď, Jene, hned ti dám prohlédnout různé části domu.

%%K t05-01-49 p
Vstupme do \VUgras{přízemí} \textsuperscript{(7)}. Zde je knoflík \VUgras{zvonku}
\textsuperscript{(11)}. Vystoupíme tři \VUgras{schody} \textsuperscript{(14)}
\VUgras{schodiště před vchodem} \textsuperscript{(9)}, překročíme \VUgras{práh}
\textsuperscript{(10)} vchodových \VUgras{dveří} \textsuperscript{(8)}, a hle, jsme u
paty \VUgras{schodiště} \textsuperscript{(13)}.

%%K t05-01-50 p
\textsc{Jan}. --- To je chodba.

%%K t05-01-51 p
\textsc{Architekt}. --- Ne, to je \VUgras{předsíň} \textsuperscript{(12)}.
\VUgras{Chodba} \textsuperscript{(a)} je na konci. Vpravo, zde je \VUgras{salon}
\textsuperscript{(b)} a \VUgras{jídelna} \textsuperscript{(c)}, naproti jídelně je
\VUgras{kuchyně} \textsuperscript{(d)} a \VUgras{pokoj pro služebné}
\textsuperscript{(e)}, potom vpředu \VUgras{pracovna} \textsuperscript{(f)}.
\VUgras{Veranda} \textsuperscript{(23)} se svými \VUgras{prosklenými dveřmi}
\textsuperscript{(25)} je blízko salonu.

%%K t05-01-52 p
Hned vystoupíme po schodišti. Přidržuj se \VUgras{zábradlí} \textsuperscript{(15)} a
počítej schody. Je jich čtrnáct. Zde je \VUgras{odpočívadlo} \textsuperscript{(m)},
jsme v prvním \VUgras{patře} \textsuperscript{(27)}.

%%K t05-tit-9 sub
\VUsaut{3.0mm}
\VUpk{10.2pt}{40}{První patro.}
\VUsaut{3.0mm}

%%K t05-01-53 p
Naproti schodišti je \VUgras{záchod} \textsuperscript{(20)} a \VUgras{toaleta}
\textsuperscript{(j)}. Vpravo krásná \VUgras{ložnice} \textsuperscript{(g)} se dvěma
okny do \VUgras{průčelí} \textsuperscript{(1)} domu. Vlevo je \VUgras{koupelna}
\textsuperscript{(1)} a \VUgras{pokoj pro hosty} \textsuperscript{(i)} s velkou
\VUgras{alkovnou} \textsuperscript{(h)}: blízko ložnice je \VUgras{dětský pokoj}
\textsuperscript{(k)}. Přiléhá k prostorné \VUgras{terase} \textsuperscript{(21)}. Pod
tou terasou je druhý záchod (20) a na zdi, zde je \VUgras{zvon}
\textsuperscript{(22)}, který zvoní k večeři.

%%K t05-01-54 p
\textsc{Jan}. --- Pojďme do \VUgras{druhého patra}.

%%K t05-01-55 p
\textsc{Architekt}. --- Druhé patro není. Pod střechou jsou jen \VUgras{podkrovní
pokojíky} \textsuperscript{(33)} a půda (34). Prohlídku jsme skončili, můžeme sejít
dolů.

%%K t05-01-56 p
\textsc{Jan}. --- Děkuji, pane, je to velmi zajímavé. Hned vypravím otci všechno, co
jsem viděl.
